% Seminar 8 handout (lecture 9): LCG full period and attack, LFSR period, matrix form,
% recovery from 2L bits (Gauss over F_2, Berlekamp–Massey).
\documentclass{ca-handout}

\seminar{8}{Генераторы псевдослучайных чисел}

\begin{document}
\makehandouttitle

\section{Линейные конгруэнтные последовательности}
$X_{n+1} = (a X_n + c) \bmod m$; \quad
$X_{n+k} \equiv a^k X_n + c\,(1 + a + \cdots + a^{k-1}) \pmod m$; период $\le m$.

\begin{reminder}[title={Теорема о максимальном периоде (Халл"--~Добелл, Кнут~A)}]
Период равен $m$ (для любого $X_0$) $\iff$
\quad 1)~$\gcd(c, m) = 1$; \quad 2)~$p \mid a - 1$ для всех простых $p \mid m$;
\quad 3)~$4 \mid a - 1$, если $4 \mid m$.
\par\smallskip
Для $m = 2^e$: $c$ нечётно и $a \equiv 1 \pmod 4$. При $c = 0$ (мультипликативный генератор)
период $= \ord_{m/d}(a)$, $d = \gcd(X_0, m)$, и не превосходит $\lambda(m)$; для простого $m$ он
максимален ($m - 1$), если $a$ "--- первообразный корень.
\end{reminder}

Если $p \mid m$, $p \nmid a - 1$, то $x^* \equiv c(1 - a)^{-1} \pmod p$ "--- неподвижная точка по
модулю $p$. Если $4 \mid m$, $a \equiv 3 \pmod 4$, то $X_{n+2} \equiv X_n \pmod 4$.

\begin{reminder}[title={Восстановление параметров по выходам $X_0, X_1, \dots$}]
$t_n = X_{n+1} - X_n$, \quad $t_{n+1} \equiv a t_n \pmod m$, \quad
$u_n = t_{n+2} t_n - t_{n+1}^2 \equiv 0 \pmod m$, \quad $m = \gcd(u_0, u_1, \dots)$;
\par
$a \equiv t_1 t_0^{-1} \pmod m$ (если $\gcd(t_0, m) = 1$), \quad $c \equiv X_1 - a X_0 \pmod m$.
\end{reminder}

\section{РСЛОС над $\GF{2}$}
$s_{n+L} = c_1 s_{n+L-1} + c_2 s_{n+L-2} + \cdots + c_L s_n$ (сложение "--- XOR).
Характеристический многочлен $f(x) = x^L + c_1 x^{L-1} + \cdots + c_L$, многочлен связей
$C(x) = 1 + c_1 x + \cdots + c_L x^L = x^L f(1/x)$; регистр невырожден, если $c_L = 1$.

\begin{reminder}[title={Матричная форма}]
$S_n = (s_n, s_{n+1}, \dots, s_{n+L-1})^T$, \quad $S_{n+1} = A S_n$, \quad $S_n = A^n S_0$, \quad
$s_n$ "--- первая координата $A^n S_0$:
\[
  A = \begin{pmatrix}
    0 & 1 & 0 & \cdots & 0\\
    \vdots & & \ddots & & \vdots\\
    0 & 0 & \cdots & 0 & 1\\
    c_L & c_{L-1} & \cdots & c_2 & c_1
  \end{pmatrix}.
\]
$A^n$ "--- бинарным алгоритмом: $n = \sum n_i 2^i$; $v \gets S_0$; для $i = 0, 1, \dots$: если
$n_i = 1$, то $v \gets A^{2^i} v$; $A^{2^{i+1}} = (A^{2^i})^2$. Проверка: $s_n = s_{n \bmod T}$.
\end{reminder}

\begin{reminder}[title={Период ($c_L = 1$, $S_0 \ne 0$)}]
$\ord_f(x)$ "--- наименьшее $e \ge 1$ с $x^e \equiv 1 \pmod f$.
\begin{itemize}
  \item Период $T \le 2^L - 1$ и $T \mid \ord_f(x)$; для приводимого $f$ он зависит от $S_0$.
  \item $f$ неприводим $\Rightarrow$ $T = \ord_f(x)$ для любого $S_0 \ne 0$, $T \mid 2^L - 1$.
  \item $T = 2^L - 1$ $\iff$ $f$ \emph{примитивен}: неприводим и
    $x^{(2^L - 1)/q} \not\equiv 1 \pmod f$ для всех простых $q \mid 2^L - 1$.
\end{itemize}
\end{reminder}

Неприводимые многочлены степени $\le 3$: $x$, $x + 1$, $x^2 + x + 1$, $x^3 + x + 1$,
$x^3 + x^2 + 1$. Многочлен степени $L \le 7$ неприводим $\iff$ он не делится ни на один из
них степени $\le L/2$.
$x + 1 \mid f \iff f(1) = 0$ (чётное число ненулевых коэффициентов).
\[
\begin{array}{c|cccc}
  L & 3 & 4 & 5 & 6\\ \hline
  2^L - 1 & 7 & 3 \cdot 5 & 31 & 3^2 \cdot 7\\
  \text{примитивные} & x^3 + x + 1 & x^4 + x + 1 & x^5 + x^2 + 1 & x^6 + x + 1
\end{array}
\]

\section{Восстановление РСЛОС}
\begin{reminder}[title={Метод Гаусса по $2L$ битам}]
\[
  \left(\begin{array}{cccc|c}
    s_0 & s_1 & \cdots & s_{L-1} & s_L\\
    s_1 & s_2 & \cdots & s_L & s_{L+1}\\
    \vdots & & & \vdots & \vdots\\
    s_{L-1} & s_L & \cdots & s_{2L-2} & s_{2L-1}
  \end{array}\right),
  \qquad\text{неизвестные } (c_L, c_{L-1}, \dots, c_1).
\]
Если линейная сложность равна $L$, решение единственно. Если ранг $r < L$, решений $2^{L - r}$;
по лемме Месси все они порождают одну и ту же последовательность.
\end{reminder}

\begin{reminder}[title={Алгоритм Берлекэмпа"--~Месси}]
$C \gets 1$, $B \gets 1$, $\ell \gets 0$, $m \gets -1$. Для $N = 0, 1, \dots, n - 1$:
$d \gets s_N + \sum_{i=1}^{\ell} c_i s_{N-i}$; если $d = 1$: $T \gets C$, $C \gets C + x^{N-m} B$,
и если $2\ell \le N$, то $\ell \gets N + 1 - \ell$, $m \gets N$, $B \gets T$.
\par Ответ: линейная сложность $\ell$, многочлен связей $C(x)$, $f(x) = x^\ell C(1/x)$.
Результат однозначен, если $n \ge 2\ell$.
\end{reminder}
\end{document}
